\begin{table}[htbp]
\centering
\small
\caption{Problema E1, Passo 5 (openai/gpt-4o-mini): viés de ordem dos exemplos.}
\label{tab:bloco2_E1_passo5_ordem_gpt4mini}
\begin{tabular}{lrrrrrr}
\toprule
 & \multicolumn{3}{c}{Acur.\ vs espec.\ (\%)} & \multicolumn{3}{c}{$\kappa$} \\ \cmidrule(lr){2-4} \cmidrule(lr){5-7}
Ordenação & k = 4 & k = 10 & k = 20 & k = 4 & k = 10 & k = 20 \\
\midrule
classe 0 primeiro & 81,5 & 80,0 & 87,7 & 0,631 & 0,600 & 0,752 \\
classe 1 primeiro & 79,0 & 84,0 & 85,1 & 0,576 & 0,682 & 0,703 \\
embaralhada & 81,7 & 86,9 & 85,6 & 0,637 & 0,738 & 0,716 \\
alternada & 81,5 & 84,8 & 86,7 & 0,629 & 0,695 & 0,735 \\
\midrule
\textbf{amplitude (máx $-$ mín)} & \textbf{2,7} & \textbf{6,9} & \textbf{2,6} &  &  &  \\
\bottomrule
\multicolumn{7}{p{0.95\linewidth}}{\footnotesize Problema E1, Passo 5: mesmos exemplos (estratégia mista) em quatro ordenações. Acur.\ vs espec.\ = acurácia (LLM vs especialista): fração dos pontos de teste em que o LLM reproduz o rótulo do especialista. $\kappa$ = Cohen's $\kappa$ (faixas de Landis e Koch, 1977). Médias das sementes.} \\
\end{tabular}
\end{table}
